The fluid model takes the form of a set of coupled PDEs governing the
species densities, temperatures, and the electric potential.  We
document a 1-D version of the model here, including the PDEs, boundary
conditions, and transport and chemistry closure models.  Extension to
2-D is necessary to model the full device, including the edge effects,
but this is not expected to be necessary for the envisioned model
validation relevant to the torch simulations.

\section{PDEs}
We require PDEs for the species densities, temperatures, and electric potential.

\subsection{Species Continuity}
The species transport equations are as follows:
%
\begin{equation}
\pp{n_{\alpha}}{t} + \pp{F_{\alpha}}{x} = \dot{\omega}_{\alpha},
\label{eqn:species}
\end{equation}
%
where $t$ is time, $x$ is the coordinate in the electrode-normal
direction, $n_{\alpha}$ denotes the number density of species
$\alpha$, $F_{\alpha}$ is the corresponding flux, and
$\dot{\omega}_{\alpha}$ is the net production rate of species $\alpha$
due to chemical reactions.  The flux is closed as follows:
%
\begin{equation*}
F_{\alpha} = z_{\alpha} \mu_{\alpha} n_{\alpha} E - D_{\alpha} \pp{n_{\alpha}}{x},
\end{equation*}
%
where $z_{\alpha}$ is the charge number, $\mu_{\alpha}$ is the
mobility, $D_{\alpha}$ is the diffusion coefficient, and $E$ is the
electric field.  The transport and chemistry closures are discussed in
\S\ref{sec:fluid-closures}.  The electric field is determined from the
potential, $E = -\pp{\phi}{x}$, which is governed by a Poisson
equation (\S\ref{sec:pdes-poisson}).

For a mixture of $N_s$ species, $N_s - 1$ species densities are
evolved according to~\eqref{eqn:species}.  The density of the dominant
background species (e.g., argon) is determined from the pressure,
which is assumed known, and the ideal gas law:
%
\begin{equation*}
p = \sum_{\alpha = 1}^{N_s} n_{\alpha} k_B T_{\alpha},
\end{equation*}
%
where $k_B$ is the Boltzmann constant and $T_{\alpha}$ is the
temperature of species $\alpha$.  As discussed in
\S\ref{sec:pdes-energy}, a two temperature model is adopted where the
heavy species share a single temperature $T_g$, such that the ideal
gas law becomes

\begin{equation*}
p = n_e k_B T_e + k_B T_g \sum_{\alpha \neq e} n_{\alpha}.
\end{equation*}



\subsection{Energy Equations} \label{sec:pdes-energy}
A two-temperature model is adopted.  We assume a common temperature
$T_g$ for the heavy species and a separate temperature for the
electrons.  The electron energy equation is given by

\begin{gather*}
\pp{}{t} \left( \frac{3}{2} k_B n_e T_e \right)
+
\pp{}{x} \left( \frac{5}{2} k_B T_e F_e - \kappa_e \pp{T_e}{x} \right)
=
\underbrace{- q_e F_e E}_{\textrm{Joule heating}}
-
\underbrace{\frac{3}{2} k_B n_e \frac{2 m_e}{m_b} \left( T_e - T_g \right) \bar{\nu}_{e,b}}_{\textrm{Elastic}}
-
\underbrace{q_e \sum_{j=1}^{N_r} \Delta E^e_j G_j}_{\textrm{Inelastic}},
\end{gather*}

where $m_e$ and $m_b$ denote the molecular mass of electron and the
background species, respectively, $\kappa_e$ is the electron thermal
conductivity, $\Delta E^e_j$ is the energy lost per electron in the
collisional process represented by reaction $j$, and $G_j$ is the rate
of progress of reaction $j$.  The thermal conductivity $\kappa_e$ and
chemistry related quantities $\Delta E_j$ and $G_j$ are described more
in \S\ref{sec:fluid-closures}.

The gas energy equation is given by

\begin{gather*}
\pp{}{t} \left( \sum_{\alpha \in h} C_{v,\alpha} n_\alpha T_g \right)
+
\pp{}{x} \left( \sum_{\alpha \in h} C_{p,\alpha} F_\alpha T_g - \sum_{\alpha \in h} \kappa_\alpha \pp{T_g}{x} \right)
=
\underbrace{C_{\textrm{inc}} \sum_{\alpha \in h} q_e z_{\alpha} F_\alpha E}_{\textrm{Joule heating}}
+
\underbrace{\frac{3}{2} k_B n_e \frac{2 m_e}{m_b} \left( T_e - T_g \right) \bar{\nu}_{e,b}}_{\textrm{Elastic}}
-
\underbrace{q_e \sum_{j=1}^{N_r} \Delta E^g_j G_j}_{\textrm{Inelastic}},
\end{gather*}

where the notation $\sum_{\alpha \in h}$ denotes the summation over
the heavy species, $C_{v,\alpha}$ and $C_{p,\alpha}$ are the specific
heat at constant volume and pressure, respectively, for species
$\alpha$.  Finally, $C_{\textrm{inc}}$ is a constant that may be used
to account for incomplete conversion of the kinetic energy gained by
the electrons into heat due to the fact that the ion mean free path
can be of the same order of the sheath thickness~\cite{}.

\subsection{Total Energy} \label{sec:pdes-energy}
Rather than solving the electron energy and heavy species thermal
energy equations, it is also possible to formulate the model in terms
of the total energy.  Specifically, let $n e_t$ denote the total
energy per unit volume:

\begin{equation*}
n e_t
=
n_e e_e + \sum_{\alpha \in h} n_{\alpha} e_{\alpha}
=
n_e \frac{3}{2} k_B T_e + \sum_{\alpha \in h} n_{\alpha} \left( \frac{3}{2} k_B T_g + \varepsilon_{\alpha} \right),
\end{equation*}

where $\varepsilon_{\alpha}$ denotes the internal energy of species
$\alpha$ (i.e., energy not associated with thermal motion).  To begin,
we assume that the heavy species are monatomic, such that the internal
energy is associated with electronic excitation.  Since these excited
states are explicitly tracked as separate species, $\varepsilon_{\alpha}$
is a given constant number, not a function requiring additional state
information or assumptions.

\todo{How will this go with more complex gases, e.g., diatomics, where
  there are more ways to store internal energy other than electronic
  excitation (vibration, rotation)?  Still track explicitly or
  introduce addtional temperatures or assume some kind of
  equilibrium?}

Further, for monatomics, $C_v = \frac{3}{2} k_B$.  As before, the
heavy species have been assumed to have the same temperature $T_g$.
Finally, we have assumed that the bulk motion of the fluid is
negligible, such that there is no contribution from the kinetic energy
associated with the bulk velocity.


This energy is governed by the following PDE:

\begin{equation}
\pp{ne_t}{t} =
\underbrace{-\,\pp{q}{x}}_{\mathrm{transport}} +
\underbrace{\sum_{\alpha} q_e z_{\alpha} F_{\alpha} E}_{\textrm{Joule heating}} +
\underbrace{\dot{Q}_{rad}}_{\textrm{Radiative heating}},
\label{eqn:totalE}
\end{equation}

where the heat flux $q$ is given by

\begin{equation}
q
=
\frac{5}{2} k_B T_e F_e - \kappa_e \pp{T_e}{x} +
\sum_{\alpha \in h} \left\{ \left( \frac{5}{2} k_B T_g + \varepsilon_{\alpha} \right) F_{\alpha} - \kappa_{\alpha} \pp{T_g}{x} \right\},
\label{eqn:totalE_flux}
\end{equation}

and the radiative heating is given by

\begin{equation}
\dot{Q}_{rad} = I dunno yet
\label{eqn:totalE_rad}
\end{equation}
\todo{Fill this in, but... we can get going assuming it is zero}


\subsubsection{Total Energy Sanity Check}
As a check on our total energy equation, we derive it from previously
introduced equations.  We have

\begin{align*}
\pp{ne_t}{t} &=
\underbrace{\pp{}{t} \left( n_e \frac{3}{2} k_B T_e \right)}_{\textrm{Electron energy}} +
\underbrace{\pp{}{t} \left( \sum_{\alpha \in h} \frac{3}{2} k_B n_\alpha T_g \right)}_{\textrm{Heavy energy}} +
\underbrace{\sum_{\alpha \in h} \pp{n_{\alpha}}{t} \varepsilon_{\alpha}}_{\textrm{Species continuity}},
\end{align*}

where the underbraces denote the previous equations that each term is
associated with.  Collecting the transport terms from these equations,
it is straightforward to show that $q$ defined
in~\eqref{eqn:totalE_flux} is consistent.  Similarly, collecting the
Joule heating terms from the electron and heavy species energy
equations, and setting $C_{\mathrm{inc}} = 1$ gives the Joule heating
term in~\eqref{eqn:totalE}.  Clearly, we may allow $C_{\mathrm{inc}}
\neq 1$ in~\eqref{eqn:totalE} by a simple modification.

From the remaining source terms we have

\begin{equation*}
-
\underbrace{\frac{3}{2} k_B n_e \frac{2 m_e}{m_b} \left( T_e - T_g \right) \bar{\nu}_{e,b}}_{\textrm{Elastic}}
-
\underbrace{q_e \sum_{j=1}^{N_r} \Delta E^e_j G_j}_{\textrm{Inelastic}},
+
\underbrace{\frac{3}{2} k_B n_e \frac{2 m_e}{m_b} \left( T_e - T_g \right) \bar{\nu}_{e,b}}_{\textrm{Elastic}}
-
\underbrace{q_e \sum_{j=1}^{N_r} \Delta E^g_j G_j}_{\textrm{Inelastic}},
+
\sum_{\alpha} \dot{\omega}_{\alpha} \varepsilon_{\alpha}.
\end{equation*}

As expected, the elastic terms simply exchange thermal energy between
the electrons and heavies and thus cancel each other.  The remaining
(inelastic) terms represent exchanges between thermal and internal
energy:

\begin{equation*}
-
\underbrace{q_e \sum_{j=1}^{N_r} \Delta E^e_j G_j}_{\textrm{Inelastic}},
-
\underbrace{q_e \sum_{j=1}^{N_r} \Delta E^g_j G_j}_{\textrm{Inelastic}},
+
\sum_{\alpha} \dot{\omega}_{\alpha} \varepsilon_{\alpha}.
\end{equation*}

To the extent that these do not balance, it must be because the
mismatch is radiated away.  For instance, we can imagine a
de-excitation event where part of the excess energy is radiated away
rather than being converted into thermal energy.  This must be
accounted for in the radiation model $\dot{Q}_{rad}$.

If we assert that $\dot{Q}_{rad} = 0$, this implies that

\begin{equation*}
\sum_{\alpha} \dot{\omega}_{\alpha} \varepsilon_{\alpha}
=
q_e \sum_{j=1}^{N_r} ( \Delta E^e_j + \Delta E^g_j) G_j
\end{equation*}

To assess this, we rewrite $\dot{\omega}_{\alpha}$ in terms of the
contribution from each reaction:

\begin{equation*}
\sum_{\alpha} \dot{\omega}_{\alpha} \varepsilon_{\alpha}
=
\sum_{j=1}^{N_r} \left\{ \sum_{\alpha} \left( \nu_{\alpha j}^{\prime \prime} - \nu_{\alpha j}^{\prime} \right) \varepsilon_{\alpha} \right\} G_j,
\end{equation*}

where $\nu_{\alpha j}^{\prime \prime}$ and $\nu_{\alpha j}^{\prime}$
are stoichiometric coefficients (see~\ref{sec:fluid-closures}).  Thus,
for $\dot{Q}_{rad} = 0$, it is sufficient to have

\begin{equation*}
\sum_{\alpha} \left( \nu_{\alpha j}^{\prime \prime} - \nu_{\alpha j}^{\prime} \right) \varepsilon_{\alpha}
=
\Delta E^e_j + \Delta E^g_j
\end{equation*}

for all reactions $j$.

\subsection{Pressure Constraint}
As noted previously, the pressure is assumed constant, which imposes
an additional constraint on the system.  There are potentially
multiple methods to implement this constraint.  Here, we assume that
the background gas (argon) is locally added or removed in order to
maintain the fixed pressure.  This results in both mass and energy
source terms in the governing PDEs.  Letting $n_b$ denote the number
density of the background species, then
%
\begin{equation*}
\pp{n_b}{t} = S,
\end{equation*}
%
where $s$ is determined by requiring constant pressure.  For a mixture
of monatomic gases,
%
\begin{equation*}
ne_t = \frac{3}{2} p + \sum_{\alpha} n_{\alpha} \varepsilon_{\alpha}.
\end{equation*}
%
Thus, since the pressure is fixed,
%
\begin{equation*}
\pp{ne_t}{t}
= \frac{3}{2} \pp{p}{t} + \sum_{\alpha} \pp{n_{\alpha}}{t} \varepsilon_{\alpha}
= \sum_{\alpha} \pp{n_{\alpha}}{t} \varepsilon_{\alpha}.
\end{equation*}
%
Using the total energy equation from before, plus a source term
imposed by addition of the background species at the rate $S$ with the
local temperature $T_g$, this imposes a constraint that determines
$S$.  Specifically,
%
\begin{equation*}
\sum_{\alpha} \pp{n_{\alpha}}{t} \varepsilon_{\alpha}
=
-\,\pp{q}{x}
+
\sum_{\alpha} q_e z_{\alpha} F_{\alpha} E
+
\dot{Q}_{rad}
+
\frac{3}{2} S k_B T_g.
\end{equation*}
%
This is an algebraic constraint that can be used to eliminate $S$.

\subsection{Electromagnetics} \label{sec:pdes-poisson}
A quasi-static model is used for the electromagnetics.  In this case,
Maxwell's equations reduce to a Poisson equation governing the
electric potential $\phi$:
%
\begin{equation*}
-\pp{^2\phi}{x^2} = \frac{\rho_c}{\epsilon_0}
\end{equation*}
%
where $\rho_c$ is the charge density and $\epsilon_0$ is the
permittivity of free space.  The charge density is given from the
species densities:
%
\begin{equation*}
\rho_c = \sum_{\alpha=1}^{N_s} z_{\alpha} q_e n_{\alpha},
\end{equation*}
%
where $q_e$ is the elementary charge (the magnitude of the charge of
an electron).
